%==========================================================================
% Section: Empirical analysis
%   Results produced by poisson_models (1).R (total claim counts) and
%   binomial_probit_models.R (extreme claim counts); see results_extreme/.
%==========================================================================
\section{Application to regional claim counts in Korea}
\label{sec.empirical}

In this section, we fit Model~\ref{model.thin} to a regional panel of
claim counts and evaluate its one-step-ahead predictions.

\subsection{Data}
\label{sec.emp.data}

The data set contains annual observations for \(I=17\) regions of
Korea over the years 2015--2024. For region \(i\) and year \(t\), the
total claim count \(N_{it}\) and the extreme claim count \(M_{it}\) are
recorded, and \(M_{it}\le N_{it}\) holds for all observations.
% TODO: state the definition of an extreme claim (e.g. the threshold) and the data source.
The covariate vector \(x_{it}\) consists of an intercept and the
logarithms of four regional characteristics: the recovery budget lagged
by one year (RecvBudget\_m1Y), the gross regional domestic product
(GRDP), the population, and the area.
% TODO: state the units of the covariates.
We use the years 2015--2023 as the training period and hold out 2024
for the evaluation of the predictions. The training period contains
\(153\) region-years, \(932\) claims, and \(200\) extreme claims. Since
\(8\) region-years have no claims, \(145\) region-years contribute to the
likelihood of the extreme claim part (Remark~\ref{rem.thin.estimation}).

Table~\ref{tab.emp.year} reports the annual totals. The total claim
count varies moderately over the years, whereas the share of extreme
claims varies considerably, from \(2.4\%\) in 2015 and \(4.1\%\) in 2021
to \(50.3\%\) in 2020. We return to this feature in
Section~\ref{sec.emp.discussion}.

\begin{table}[htbp]
\centering
\caption{Annual totals of the total and extreme claim counts over the
17 regions.}
\label{tab.emp.year}
\small
\begin{tabular}{lrrrrrrrrrr}
\toprule
Year & 2015 & 2016 & 2017 & 2018 & 2019 & 2020 & 2021 & 2022 & 2023 & 2024\\
\midrule
\(\sum_i N_{it}\) & 42 & 74 & 82 & 121 & 101 & 151 & 123 & 99 & 139 & 151\\
\(\sum_i M_{it}\) & 1 & 20 & 8 & 9 & 12 & 76 & 5 & 26 & 43 & 32\\
Share (\%) & 2.4 & 27.0 & 9.8 & 7.4 & 11.9 & 50.3 & 4.1 & 26.3 & 30.9 & 21.2\\
\bottomrule
\end{tabular}
\end{table}

\subsection{Estimation}
\label{sec.emp.estimation}

By Remark~\ref{rem.thin.estimation}, the maximum likelihood estimators
of \((\alpha,r)\) and \((\beta,\sigma)\) are obtained separately from
\(\ell^{(1)}\) and \(\ell^{(2)}\). The log-likelihood \(\ell^{(1)}\) is
available in closed form by \eqref{eq.thin.L1}. The log-likelihood
\(\ell^{(2)}\) is evaluated by the one-dimensional integral in
Remark~\ref{rem.thin.computation} using Gauss--Hermite quadrature with
\(100\) nodes; increasing the number of nodes to \(200\) changes
\(\ell^{(2)}\) by less than \(10^{-6}\), and adaptive numerical
integration agrees to within \(2\times10^{-4}\). Both log-likelihoods
are maximized by a quasi-Newton method with analytic gradients followed
by Newton iterations, and the standard errors are obtained from the
observed information matrix. The standard errors of the random effect
variances \(1/r\) and \(\sigma^2\) are obtained by the delta method.

As a benchmark, we fit the corresponding models without random effects,
namely the Poisson generalized linear model (GLM) for \(N_{it}\) and the
binomial probit GLM for \(M_{it}\) given \(N_{it}\). These correspond to
Model~\ref{model.thin} with \(\Theta_i^{(1)}\equiv1\) and
\(\Theta_i^{(2)}\equiv0\), so that the predictions are
\(\lambda_{i,t+1}\) and \(\Phi(\eta_{i,t+1})\), without experience
rating. We consider two specifications of \(x_{it}\): Specification A
uses all four covariates, and Specification B omits the recovery budget.
The random effect models are referred to as the generalized linear mixed
models (GLMMs).

Table~\ref{tab.emp.N} reports the estimates of the total claim count
part. The gamma random effect is supported by the likelihood ratio test,
whose \(p\)-value is computed from the boundary-corrected
\(\frac12\chi^2_1\) distribution, and by both information criteria. The
recovery budget is not significant in either model, and removing it
lowers both information criteria.

\begin{table}[htbp]
\centering
\caption{Estimates of the total claim count part \((\alpha,r)\), with
standard errors in parentheses. GLM: Poisson GLM; GLMM: Poisson-gamma
GLMM.}
\label{tab.emp.N}
\small
\begin{tabular}{lrrrr}
\toprule
& \multicolumn{2}{c}{Specification A} & \multicolumn{2}{c}{Specification B}\\
\cmidrule(lr){2-3}\cmidrule(lr){4-5}
& GLM & GLMM & GLM & GLMM\\
\midrule
Intercept & \(-3.192\) (1.024) & \(-0.096\) (2.492) & \(-2.559\) (0.737) & \(0.399\) (1.939)\\
log(RecvBudget\_m1Y) & \(0.060\) (0.067) & \(0.033\) (0.111) & -- & --\\
log(GRDP) & \(0.313\) (0.130) & \(1.120\) (0.432) & \(0.327\) (0.129) & \(1.179\) (0.390)\\
log(Population) & \(-0.117\) (0.134) & \(-0.910\) (0.415) & \(-0.115\) (0.133) & \(-0.958\) (0.389)\\
log(Area) & \(0.247\) (0.047) & \(0.230\) (0.088) & \(0.283\) (0.026) & \(0.247\) (0.068)\\
\(1/r\) & -- & \(0.091\) (0.066) & -- & \(0.100\) (0.064)\\
\midrule
Log-likelihood & \(-388.78\) & \(-382.69\) & \(-389.19\) & \(-382.73\)\\
AIC & \(787.55\) & \(777.38\) & \(786.38\) & \(775.47\)\\
BIC & \(802.71\) & \(795.56\) & \(798.50\) & \(790.62\)\\
LR test (\(p\)-value) & \multicolumn{2}{c}{\(12.18\) (\(0.0002\))} & \multicolumn{2}{c}{\(12.91\) (\(0.0002\))}\\
\bottomrule
\end{tabular}
\end{table}

Table~\ref{tab.emp.M} reports the estimates of the extreme claim part.
As for the total claim count, the Gaussian random effect is supported
by the likelihood ratio test in both specifications, with
\(\widehat\sigma^2=0.117\) in Specification A and
\(\widehat\sigma^2=0.058\) in Specification B, and the GLMM attains the
lower AIC and BIC. The recovery budget is again not significant, and
Specification B has the lower AIC and BIC among the GLMMs. In both
parts, the coefficients of log(GRDP) and log(Population) have opposite
signs and change considerably once the random effect is introduced.
These two covariates are strongly correlated across regions, and
population varies little within a region over time, so that only their
contrast, a GRDP per capita effect, is well identified. The coefficient
of log(Area) is the most stable across models; a larger area is
associated with a higher frequency and a higher share of extreme claims.
These associations are descriptive and should not be interpreted as
causal effects.

\begin{table}[htbp]
\centering
\caption{Estimates of the extreme claim part \((\beta,\sigma)\), with
standard errors in parentheses.}
\label{tab.emp.M}
\small
\begin{tabular}{lrrrr}
\toprule
& \multicolumn{2}{c}{Specification A} & \multicolumn{2}{c}{Specification B}\\
\cmidrule(lr){2-3}\cmidrule(lr){4-5}
& Probit GLM & Probit GLMM & Probit GLM & Probit GLMM\\
\midrule
Intercept & \(-0.220\) (1.735) & \(3.216\) (3.269) & \(0.246\) (1.105) & \(-0.098\) (1.808)\\
log(RecvBudget\_m1Y) & \(0.039\) (0.111) & \(-0.294\) (0.206) & -- & --\\
log(GRDP) & \(0.521\) (0.203) & \(1.033\) (0.448) & \(0.531\) (0.201) & \(0.733\) (0.305)\\
log(Population) & \(-0.650\) (0.209) & \(-0.968\) (0.420) & \(-0.652\) (0.208) & \(-0.779\) (0.316)\\
log(Area) & \(0.261\) (0.077) & \(0.444\) (0.148) & \(0.284\) (0.044) & \(0.269\) (0.063)\\
\(\sigma^2\) & -- & \(0.117\) (0.088) & -- & \(0.058\) (0.041)\\
\midrule
Log-likelihood & \(-310.05\) & \(-305.34\) & \(-310.12\) & \(-306.63\)\\
AIC & \(630.10\) & \(622.67\) & \(628.23\) & \(623.25\)\\
BIC & \(644.98\) & \(640.53\) & \(640.14\) & \(638.14\)\\
LR test (\(p\)-value) & \multicolumn{2}{c}{\(9.43\) (\(0.0011\))} & \multicolumn{2}{c}{\(6.98\) (\(0.0041\))}\\
Pearson \(\chi^2\)/df & \(4.32\) & -- & \(4.27\) & --\\
\bottomrule
\end{tabular}
\end{table}

\subsection{One-step-ahead prediction of 2024}
\label{sec.emp.prediction}

Using the training years as the history \(\mathcal H_{i,2023}\), we
compute the predictive probability \(\widehat p_{i,2024}\) of an extreme
claim from part iii of Theorem~\ref{thm.thin.pred}, and predict
\(M_{i,2024}\) in two ways:
\begin{itemize}
  \item[(a)] given the realized total claim count, by
  \(N_{i,2024}\,\widehat p_{i,2024}\) (part iv of
  Theorem~\ref{thm.thin.pred});
  \item[(b)] without it, by
  \(\E{N_{i,2024}\mid\mathcal H_{i,2023}}\,\widehat p_{i,2024}\)
  (part v of Theorem~\ref{thm.thin.pred}), where the first factor is
  given by part ii of Theorem~\ref{thm.thin.pred} with the estimates in
  Table~\ref{tab.emp.N}.
\end{itemize}
For the GLM benchmark, \(\widehat p_{i,2024}\) and
\(\E{N_{i,2024}\mid\mathcal H_{i,2023}}\) are replaced by
\(\Phi(\eta_{i,2024})\) and \(\lambda_{i,2024}\), respectively, with
the GLM estimates. The predictions are evaluated by the root mean
squared error (RMSE), the mean absolute error (MAE), and the deviance,
which is the binomial deviance in case (a) and the Poisson deviance in
case (b).

Table~\ref{tab.emp.crit} summarizes the results. In contrast to the
in-sample comparison, the GLMM does not improve the out-of-sample
predictions of the extreme claim counts: the GLM attains slightly lower
values of all three criteria in both cases and both specifications. In
case (b), the GLMM predicts a total of \(33.8\) and \(34.6\) extreme
claims, which is closer to the realized total of \(32\) than the GLM
totals of \(28.4\) and \(27.8\); the improvement in the total stems
from the total claim count part, for which the GLMM predicts the
realized total of \(151\) claims better than the GLM (\(133.4\) and
\(133.9\) against \(113.0\) and \(111.6\)).

\begin{table}[htbp]
\centering
\caption{Prediction criteria for the extreme claim counts in 2024
(\(17\) regions, realized total \(\sum_iM_{i,2024}=32\)).}
\label{tab.emp.crit}
\small
\begin{tabular}{llrrrr}
\toprule
& & \multicolumn{2}{c}{Specification A} & \multicolumn{2}{c}{Specification B}\\
\cmidrule(lr){3-4}\cmidrule(lr){5-6}
& Criterion & GLM & GLMM & GLM & GLMM\\
\midrule
(a) given \(N_{i,2024}\) & RMSE & \(2.341\) & \(2.524\) & \(2.323\) & \(2.569\)\\
& MAE & \(1.556\) & \(1.762\) & \(1.534\) & \(1.786\)\\
& Binomial deviance & \(37.10\) & \(43.30\) & \(36.66\) & \(43.83\)\\
& Predicted total & \(39.7\) & \(41.1\) & \(39.4\) & \(41.7\)\\
\midrule
(b) unconditional & RMSE & \(1.655\) & \(1.780\) & \(1.634\) & \(1.790\)\\
& MAE & \(1.170\) & \(1.304\) & \(1.157\) & \(1.326\)\\
& Poisson deviance & \(21.74\) & \(24.09\) & \(21.14\) & \(24.50\)\\
& Predicted total & \(28.4\) & \(33.8\) & \(27.8\) & \(34.6\)\\
\bottomrule
\end{tabular}
\end{table}

Table~\ref{tab.emp.region} reports the region-level results for
Specification B, ordered by the posterior mean of
\(\Theta_i^{(2)}\). The posterior means range from \(-0.30\) for
Chungnam to \(0.28\) for Gwangju, and the credibility adjustment moves
\(\widehat p_{i,2024}\) away from the a priori probability
\(\Phi\left(\eta_{i,2024}/\sqrt{1+\widehat\sigma^2}\right)\) by up to
\(11\) percentage points. The largest prediction errors of the GLMM
occur in regions whose share of extreme claims in 2024 departs from
their own history. For example, Gyeonggi and Chungnam had \(16\) of
\(119\) and \(14\) of \(70\) claims classified as extreme during the
training period, which lowers their predictive probabilities, but
\(7\) of \(13\) and \(7\) of \(12\) claims in 2024. Conversely, Jeonnam
and Gyeongnam had high shares during the training period but only
\(4\) of \(17\) and \(2\) of \(19\) in 2024.

\begin{table}[htbp]
\centering
\caption{Region-level results for Specification B: training totals,
posterior mean of \(\Theta_i^{(2)}\), a priori and predictive
probabilities of an extreme claim, and realized and predicted extreme
claim counts in 2024 given \(N_{i,2024}\) (case (a)).}
\label{tab.emp.region}
\footnotesize
\setlength{\tabcolsep}{4pt}
\begin{tabular}{lrrrrrrrrr}
\toprule
& \multicolumn{2}{c}{2015--2023} & & \multicolumn{3}{c}{Probability of extreme claim} & \multicolumn{3}{c}{2024}\\
\cmidrule(lr){2-3}\cmidrule(lr){5-7}\cmidrule(lr){8-10}
Region & \(\sum_tN_{it}\) & \(\sum_tM_{it}\) & \(\E{\Theta_i^{(2)}\mid\mathcal H}\) & A priori & GLM & GLMM & \(N_{i,2024}\) & \(M_{i,2024}\) & GLMM pred.\\
\midrule
Gwangju & 28 & 5 & \(0.278\) & 0.079 & 0.076 & 0.125 & 5 & 0 & 0.62\\
Jeonnam & 58 & 25 & \(0.214\) & 0.383 & 0.367 & 0.463 & 17 & 4 & 7.87\\
Busan & 46 & 5 & \(0.165\) & 0.085 & 0.075 & 0.110 & 5 & 0 & 0.55\\
Jeonbuk & 68 & 18 & \(0.162\) & 0.241 & 0.249 & 0.289 & 8 & 4 & 2.31\\
Seoul & 35 & 6 & \(0.145\) & 0.132 & 0.086 & 0.161 & 2 & 0 & 0.32\\
Gyeongnam & 64 & 19 & \(0.129\) & 0.301 & 0.284 & 0.344 & 19 & 2 & 6.53\\
Gyeongbuk & 94 & 38 & \(0.129\) & 0.393 & 0.376 & 0.440 & 9 & 3 & 3.96\\
Chungbuk & 70 & 22 & \(0.067\) & 0.333 & 0.317 & 0.354 & 10 & 3 & 3.54\\
Daejeon & 22 & 1 & \(-0.030\) & 0.082 & 0.079 & 0.077 & 6 & 1 & 0.46\\
Gangwon & 82 & 21 & \(-0.035\) & 0.333 & 0.345 & 0.317 & 22 & 1 & 6.97\\
Daegu & 22 & 1 & \(-0.036\) & 0.097 & 0.097 & 0.090 & 6 & 0 & 0.54\\
Sejong & 20 & 1 & \(-0.067\) & 0.100 & 0.110 & 0.088 & 2 & 0 & 0.18\\
Ulsan & 32 & 4 & \(-0.144\) & 0.260 & 0.220 & 0.214 & 6 & 0 & 1.28\\
Jeju & 45 & 3 & \(-0.146\) & 0.155 & 0.170 & 0.121 & 0 & 0 & 0.00\\
Incheon & 57 & 1 & \(-0.252\) & 0.115 & 0.101 & 0.072 & 9 & 0 & 0.65\\
Gyeonggi & 119 & 16 & \(-0.282\) & 0.280 & 0.227 & 0.191 & 13 & 7 & 2.48\\
Chungnam & 70 & 14 & \(-0.298\) & 0.392 & 0.354 & 0.283 & 12 & 7 & 3.39\\
\bottomrule
\end{tabular}
\end{table}

\subsection{Discussion}
\label{sec.emp.discussion}

The results indicate that the static random effects of
Model~\ref{model.thin} capture persistent differences between regions,
as shown by the likelihood ratio tests, but that the share of extreme
claims is dominated by shocks that are common within a year. Two
diagnostics support this interpretation. First, the Pearson statistic of
the probit GLM in Table~\ref{tab.emp.M} is about \(4.3\) times its
degrees of freedom, so the extreme claim counts are strongly
overdispersed relative to the binomial distribution. Second, adding
year fixed effects to the probit GLM of Specification B reduces the
deviance by \(156.9\) on \(8\) degrees of freedom, whereas adding region
fixed effects on top of them reduces it by a further \(41.0\) on \(16\)
degrees of freedom. Even with both sets of fixed effects, the Pearson
statistic remains about \(3.9\) times its degrees of freedom. This
suggests that the claims within a region-year are not conditionally
independent, as would be the case when several claims arise from the
same event. Under such shocks, the experience of a region over the
training period is only weakly informative about its share of extreme
claims in the next year, which explains why the credibility adjustment
of \(\widehat p_{i,2024}\) does not improve the predictions.

These findings are subject to two caveats. The evaluation is based on a
single hold-out year with \(17\) regions and \(32\) extreme claims, so
the differences in Table~\ref{tab.emp.crit} are small relative to their
sampling variability; a rolling-origin evaluation over several years
would give a more reliable comparison. Moreover, the conclusions concern
the static specification. Natural extensions of
Model~\ref{model.thin} are a common year effect, a random effect at the
region-year level to absorb the dependence among claims from the same
event, and dynamic random effects that discount older experience.
